\documentclass{zkdl-presentation-template}

\title[Lookup Checks]{\textbf{Lookup Checks. Plookup. Logup}}
\author{Blockstream Research}
\date{\today}
\homepage{research.blockstream.com}
\github{BlockstreamResearch/zk-book}

\begin{document}
    \maketitleframe

    \begin{frame}{Plan}
      \tableofcontents
    \end{frame}

	\section{Introduction}

    \begin{frame}{Motivation}
        Suppose you want to implement the AES-128 or SHA-256 using arithmetical
        circuits (for instance, for national passport verification).\pause

        As the part of such algorithm, assume given
        $\boldsymbol{a},\boldsymbol{b},\boldsymbol{c} \in \{0,1\}^n$, you need to implement XORing:
        \begin{equation*}
            \boldsymbol{c} = \boldsymbol{a} \oplus \boldsymbol{b}
        \end{equation*}

        \pause\begin{block}{Arithmetical Circuit, informally}
            \begin{itemize}
                \item Bit-decompose $\boldsymbol{a}$ to $\{a_i\}_{i \in [n]}$ and check $a_i(1-a_i)=0, i\in [n]$.\pause
                \item Bit-decompose $\boldsymbol{b}$ to $\{b_i\}_{i \in [n]}$ and check $b_i(1-b_i)=0, i\in [n]$.\pause
                \item Verify $c_i = a_i+b_i-2a_ib_i$ (which is precisely
                $c_i=a_i \oplus b_i$).\pause
            \end{itemize}

            \textbf{In Total:} $3n$ constraints for \textit{one XORing
            operation}. E.g., for 8-bit XORing we would have approximately 24
            constraints per operation.
        \end{block}
    \end{frame}

    \begin{frame}{Motivation (cont.)}
        Suppose you need to compute \textbf{10k} 8-bit XOR operations during
        hashing (that corresponds to $\approx \text{four}$ SHA256 512-bit
        blocks). 
        \begin{equation*}
            \text{Total constraints} = \text{10k} \times 24 = 240k \; \text{constraints}
        \end{equation*}

        \pause\begin{alertblock}{Lookup}
            With lookup, you pay $2^{16}$ constraints to commit to the lookup
            table, and then you use only $1$ constraint per XOR. Thus in total
            you would have $2^{16} + 10k \approx 75k$ constraints. So we get
            $3.2\times$ boost!\pause
        \end{alertblock}

        More generally, lookup check allows to reduce complexity from
        $O(mn)$ (where $m$ is the cost of each operation and $n$ is
        the number of operations) to $O(n+d)$, where $d$ is the lookup
        size.
    \end{frame}

    \begin{frame}{Application to Bionetta}
        \begin{figure}
            \centering
            \includegraphics[width=0.15\textwidth]{images/lecture_17/logo.png}
        \end{figure}

        In \textit{Bionetta}, lookup checks allow to reduce the number of
        constraints $O(n)$ (where $n$ is the circuit size in Groth16)
        to sublinear $O(n/\log n)$. \pause

        More specifically, implementing $\text{ReLU}(x) = \max\{0,x\}$ naively
        requires roughly $O(b)$ constraints where $b = \log_2\#\mathbb{F}$.\pause
        
        In \textit{UltraGroth}, by splitting the integer into $w$-bit limbs, one
        can reduce the complexity per ReLU down to $b/w$ with the lookup
        commitment cost of $2^w$ constraints. Thus, in total, one gets
        $O(2^w + \ell b/w)$ constraints with $\ell$ being the number
        of ReLUs.\pause

        For optimal value of $w$, this reduces to $O(\ell b/\log
        (\ell b))$.
    \end{frame}

    \begin{frame}{Formalizing the problem}
        \begin{definition}
            Lookup check consists in proving/verifying that $\{z_i\}_{i \in [n]}
            \subseteq \{t_j\}_{j \in [d]}$. We will denote multisets as $\vec{z}
            := \{z_i\}_{i \in [n]}$ and $\vec{t} = \{t_j\}_{j \in [d]}$ and
            write inclusion check $\vec{z} \subseteq \vec{t}$ for short.\pause
        \end{definition}

        \begin{example}
            \begin{itemize}
                \item $\vec{t}=\{0,1\}$. $\vec{z} \subseteq \vec{t}$ means ``all $z_i$'s are binary''.\pause
                \item $\vec{z}=\{10,6,7,1,1,6,10,7,1\} \subseteq \{1,6,7,10\} = \vec{t}$. On the other hand, we 
                have $\vec{z}^* := \{1,6,10,5\} \not\subseteq \vec{t}$.\pause
                \item For $\vec{t} = \{t_j\}_{j \in [2^w]}$ with $t_j=j$,
                condition ``$\vec{z} \subseteq \vec{t}$'' means ``every element
                $z_i$ is a $w$-bit integer''.
            \end{itemize}
        \end{example}
    \end{frame}

    \begin{frame}{Working with tuples}
        In particular, XOR example can also be reduced to this check. 
        \begin{itemize}
            \item First initialize all tuples $\vec{t} := \{(a_i,b_i,c_i)\}_{i
            \in [2^{2n}]}$ such that $a_i \oplus b_i = c_i$ (perceive each
            $a_i,b_i,c_i$ as a field element from $\mathbb{F}$).\pause
            \item We need to check whether $\vec{z} := \{(x_i,y_i,w_i)\}_{i \in
            [m]} \subseteq \vec{t}$. The verifier samples the challenge $\alpha
            \leftarrowS \mathbb{F}$, and lookup table is perceived as
            $\vec{t}_{\alpha} = \{a_i+\alpha b_i+\alpha^2 c_i\}_{i \in [2^{2n}]}$,
            witness as $\vec{z}_{\alpha} := \{x_i+\alpha y_i + \alpha^2 w_i\}_{i
            \in [m]}$, and then the check is $\vec{z}_{\alpha} \subseteq
            \vec{t}_{\alpha}$ as usual.\pause
        \end{itemize}

        \begin{remarkblock}{Conclusion}
            Lookup checks are cool, so let us study them!
        \end{remarkblock}
    \end{frame}

    \section{Plookup}

    \begin{frame}{plookup}
        \begin{definition}
            Plookup is the Poly-IOP-based lookup check protocol that uses rather
            \textit{exotic} multiset equality check. Why I call it
            \textit{exotic} you will see in just a moment.\pause
        \end{definition}

        \textbf{Reminder.} When constructing $\mathcal{P}$lon$\mathcal{K}$, we considered the
        so-called \textit{permutation check}, which checked whether
        multisets $\{a_i\}_{i \in [n]}$ and $\{b_i\}_{i \in [n]}$ are equal.\pause\hspace{1.0px} This was done by running 
        the grand product:
        \begin{equation*}
            \prod_{i \in [n]} (\gamma+a_i) = \prod_{i \in [n]}(\gamma+b_i), \quad \gamma \leftarrowS \mathbb{F}
        \end{equation*}

        \vspace{-15px}

        \pause\textit{Can we use it for checking $\{z_i\}_{i\in [n]} \subseteq
        \{t_j\}_{j \in [d]}$?}\pause\hspace{1.0px} No, consider:
        \begin{equation*}
            Z(\gamma) = \prod_{i \in [n]} (\gamma + z_i), \quad T(\gamma) = \prod_{j \in [d]} (\gamma + t_j)
        \end{equation*}

        \vspace{-15px}

        We can merely state that roots of $Z(\gamma)$ and $T(\gamma)$ are the
        same.
    \end{frame}

    \begin{frame}{plookup: problematic solution}
        \begin{alertblock}{Core problem}
            How to reduce lookup check to the multiset equality check?\pause
        \end{alertblock}

        \begin{definition}
            Given $\vec{s} = \{s_i\}_{i \in [n]}$, denote by $\partial\vec{s}$
            the \textbf{difference set} $\{s_{i+1}-s_i\}_{i \in [n-1]}$ without
            zero elements. For example, $\partial\{1,1,2,5,5\} = \{1,3\}$.\pause
        \end{definition}

        \textbf{Attempt \#1.} Suppose we sort witness $\vec{z}$ and get
        $\vec{s}$ as a result. We might reasonably expect that $\partial\vec{s}
        = \partial\vec{t}$.\pause
        \begin{example}
            Suppose $\vec{t}=\{1,4,8\}$ and $\vec{s} = \{1,1,4,8,8,8\}$. Notice
            that $\vec{s} \subseteq \vec{t}$ and moreover $\partial\vec{t} =
            \partial\vec{s} = \{3,4\}$.\pause
        \end{example}

        \textcolor{oc-red-8}{\textbf{Problem.}} Converse is false. Consider $\vec{s}^*:=\{1,5,5,5,8,8\}$.
    \end{frame}

    \begin{frame}{plookup: solution \#1}
        \textbf{Attempt \#2.} Construct $\vec{s}$ --- sorted concatenation
        $(\vec{z},\vec{t})$. First assert permutation check of $\vec{s}$ and
        $(\vec{z},\vec{t})$. Then, assert that $\partial\vec{s} =
        \partial\vec{t}$.\pause
        
        \begin{proposition}
            This is necessary and sufficient condition for $\vec{z} \subseteq
            \vec{t}$.\pause
        \end{proposition}

        \begin{example}
            Suppose $\vec{t} = \{1,4,8\}$, $\vec{z} = \{1,1,4,8,8,8\}$, and
            $\vec{z}^*=\{1,5,5,5,8,8\}$. For $\vec{t}$ and $\vec{z}$ we see that
            $\vec{s} = \{1,1,1,4,4,8,8,8,8\}$ and clearly $\partial\vec{s} =
            \{3,4\} = \partial\vec{t}$. As for $\vec{z}^*$, the sorted
            concatenation is $\vec{s}^*=\{1,1,4,5,5,5,8,8,8\}$ and so
            $\partial\vec{s}^*=\{1,3\} \neq \partial\vec{t}$.\pause
        \end{example}

        Thus our protocol might simply be running two permutation checks: (a) on
        $\vec{s}$ and $(\vec{z},\vec{t})$, and (b) on $\partial\vec{s}$ and
        $\partial\vec{t}$. However, that's still \textbf{two} grand products. We
        can reduce this to a single grand product.
    \end{frame}

    \begin{frame}{plookup: solution \#2}
        \textbf{Attempt \#3.} Why do we need this strange difference set $\partial\vec{s}$? Let us 
        make it more general!\pause
        \begin{definition}
            The \textbf{randomized difference set} of $\vec{s}=\{s_i\}_{i \in
            [n]}$ for randomness $\beta$, denoted as $\partial_{\beta}\vec{s}$,
            is defined as $\{s_i+\beta s_{i+1}\}_{i \in [n-1]}$.\pause
        \end{definition} 

        \begin{proposition}
            Necessary and sufficient condition for lookup check is
            $\partial_{\beta}\vec{s} = ((1+\beta)\vec{z},
            \partial_{\beta}\vec{t})$ for randomly chosen $\beta \leftarrowS
            \mathbb{F}^{\times}$.\pause
        \end{proposition}

        \textbf{Intuition.} Note that for two same consecutive integers
        $s_i,s_{i+1}$, instead of zero, one gets $(1+\beta)s_i$. Other than
        that, we also have elements of form $t_i+\beta t_{i+1}$, which obviously
        form $\partial_{\beta}\vec{t}$. Thus, concatenating all elements of form
        $(1+\beta)s_i$ and $\partial_{\beta}\vec{t}$ gives
        $\partial_{\beta}\vec{s}$.
    \end{frame}

    \begin{frame}{Illustration}
        Suppose we, again, have the following witness and table:
        \begin{equation*}
            \vec{t} = \{1,4,8\}, \quad \vec{z} = \{1,1,4,8,8,8\}
        \end{equation*}

        \pause The sorted array is $\vec{s} = \{1,1,1,4,4,8,8,8,8\}$. Now sample 
        random $\beta \leftarrowS \mathbb{F}^{\times}$. Then we have:
        \begin{equation*}
            \partial_{\beta}\vec{s} = \{1+\beta,1+\beta,1+4\beta,(1+\beta)4,4+8\beta,(1+\beta)8,(1+\beta)8,(1+\beta)8\}
        \end{equation*}

        \pause On the other hand, we also have:
        \begin{gather*}
            (1+\beta)\vec{z} = \{1+\beta,1+\beta,(1+\beta)4,(1+\beta)8,(1+\beta)8,(1+\beta)8\} \\
            \partial_{\beta}\vec{t} = \{1+4\beta, 4+8\beta\}
        \end{gather*}
        \pause Clearly,
        $\partial_{\beta}\vec{s}=((1+\beta)\vec{z},\partial_{\beta}\vec{t})$.
    \end{frame}

    \begin{frame}{Protocol Specifics}
        Now, given $\boldsymbol{t} \in \mathbb{F}^d$, $\boldsymbol{z} \in \mathbb{F}^n$,
        and $\boldsymbol{s} \in \mathbb{F}^{n+d}$, define bi-variate polynomials $Z$ and
        $T$ as follows:
        \begin{gather*}
            Z(\beta,\gamma) \triangleq (1+\beta)^n \prod_{i \in [n]}(\gamma+z_i)\prod_{i \in [d-1]}(\gamma(1+\beta)+t_i+\beta t_{i+1}) \\
            T(\beta,\gamma) \triangleq \prod_{i \in [n+d-1]} (\gamma(1+\beta)+s_i+\beta s_{i+1})
        \end{gather*}

        \pause\begin{proposition}
            $Z \equiv T$ if and only if $\boldsymbol{z} \subseteq \boldsymbol{t}$ and
            $\boldsymbol{s}$ is $(\boldsymbol{z},\boldsymbol{t})$ sorted by
            $\boldsymbol{t}$.\pause
        \end{proposition}

        Similarly to permutation check equation
        $\prod_{j}(\gamma+a_j)=\prod_j(\gamma+b_j)$, the rest of the protocol is
        done to ensure that $Z=T$ in the $\mathcal{P}$lon$\mathcal{K}$ish manner. See lecture notes for 
        concrete details.
    \end{frame}

    \section{Logup}

    \begin{frame}{Derivatives}
        Similarly to calculus, we can define derivative operations over polynomials 
        and rational function over arbitrary fields.\pause
        \begin{definition}
            Given a polynomial $q(T) := \sum_{j=0}^d q_jT^j$ over $\mathbb{F}[T]$, the \textbf{formal
            derivative}, denoted by $q'(T)$, is given by
            $\sum_{j=1}^{d}jq_jT^{j-1}$.\pause
        \end{definition}
        \begin{definition}
            For a function $q(T)/r(T)$ from rational function field
            $\mathbb{F}(T)$, the \textbf{formal derivative} is given by
            \begin{equation*}
                \left(\frac{q(T)}{r(T)}\right)' = \frac{q'(T)r(T) - q(T)r'(T)}{r(T)^2}
            \end{equation*}
        \end{definition}
    \end{frame}

    \begin{frame}{Logarithmic Derivative}
        \begin{definition}
            The \textbf{logarithmic derivative} of $q(T) \in \mathbb{F}[T]$ is given by the 
            rational function $\mathsf{logD}[q] = q'(T)/q(T)$.\pause
        \end{definition}
        \begin{remarkblock}{Motivation}
            What is the derivative for $\log f(x)$ given $f: \mathbb{R} \to
            \mathbb{R}$?\pause \hspace{1px} By the chain rule, $f'(x)/f(x)$, which is given by
            the definition above.\pause
        \end{remarkblock}

        \begin{lemma}
            $\mathsf{logD}[qr]=\mathsf{logD}[q] + \mathsf{logD}[r]$.\pause
        \end{lemma}

        Consequently, we can infer:
        \begin{equation*}
            \mathsf{logD}\left[\prod_{j \in [n]} (T+a_j)\right] = \sum_{j \in [n]} \mathsf{logD}\left[T+a_j\right] = \sum_{j \in [n]} \frac{1}{T+a_j}
        \end{equation*}
    \end{frame}

    \begin{frame}{Finally something useful}
        \begin{lemma}[On fractional permutation check]
            Let $\vec{a} = \{a_i\}_{i \in [n]}$ and $\vec{b} = \{b_i\}_{i \in [n]}$ be
            two sequences of elements from $\mathbb{F}$. To verify with overwhelming
            probability whether two multisets are equal, it suffices to check
            \begin{equation*}
                \sum_{j \in [n]} \frac{1}{\gamma + a_j} = \sum_{j \in [n]} \frac{1}{\gamma+b_j}
            \end{equation*}
            for randomly chosen $\gamma \leftarrowS \mathbb{F}$.\pause
        \end{lemma}

        \textbf{Intuition.} Take logarithmic derivative from both sides of a
        permutation check equation $\prod_{j \in [n]} (\gamma+a_j) = \prod_{j \in [n]}
        (\gamma+b_j)$ and you are done (see prev. slide). Other direction is a
        bit trickier to prove, but still trivial enough.
    \end{frame}

    \begin{frame}{Central Equation}
        \begin{lemma}[On fractional lookup check]
            Given two sequences of elements $\{t_i\}_{i \in [d]}$ and $\{z_i\}_{i \in
            [n]}$, the inclusion check $\{z_i\}_{i \in [n]} \subseteq \{t_i\}_{i \in
            [d]}$ is satisfied if and only if there exist the set of multiplicities
            $\{\mu_i\}_{i \in [d]}$ where $\mu_i = \#\{j \in [n]: z_j = t_i\}$ such that:
            \begin{equation*}
                \sum_{i \in [n]} \frac{1}{T+z_i} = \sum_{i \in [d]} \frac{\mu_i}{T+t_i}
            \end{equation*}
        \end{lemma}

        \pause In particular, checking such equality at random point from $\mathbb{F}$
        results in the soundness error of up to $(n+d)/\#\mathbb{F}$, which becomes
        negligible for fairly large $\#\mathbb{F}$.\pause

        \begin{remarkblock}{Note}
            This is the central equation used in a large number of studies on its own. What follows 
            is just one variation of how to apply SumCheck using this equation.
        \end{remarkblock}
    \end{frame}

    \begin{frame}{Bringing Sum-Check}
        Suppose the lookup table size is $d=2^v$ for some $v$. Then, the table
        $\vec{t} = \{t_j\}_{j \in [d]}$ can be viewed as a function $t:
        \{0,1\}^v \to \mathbb{F}$.\pause

        \textcolor{oc-red-8}{\textbf{Problem:}} attempting to define $\vec{z} = \{z_i\}_{i \in [n]}$
        in a similar manner fails since in practice $n>d$, so the function
        $\{0,1\}^v \to \mathbb{F}$ has a too-small domain. Instead, split $\vec{z}$ into 
        $m=\lceil n/d \rceil$ f-ns $z_1,\dots,z_m: \{0,1\}^v \to \mathbb{F}$.\pause

        \begin{remarkblock}{Note}
            In other words, we reduced the problem of $\{z_i\}_{i \in [n]}
            \subseteq \{t_j\}_{j \in [d]}$ to $\bigcup_{i \in
            [m]}\{z_i(\mathbf{x})\}_{\mathbf{x} \in \{0,1\}^v} \subseteq
            \{t(\mathbf{x})\}_{\mathbf{x} \in \{0,1\}^v}$.\pause
        \end{remarkblock}

        Sum-Check equation is as follows:
        \begin{equation*}
            \boxed{\sum_{\mathbf{x} \in \{0,1\}^v} \sum_{i \in [m]} \frac{1}{\gamma + z_i(\mathbf{x})} = \sum_{\mathbf{x} \in \{0,1\}^v} \frac{\mu(\mathbf{x})}{\gamma+t(\mathbf{x})}},
        \end{equation*}
        where $\mu(\mathbf{x}) = \sum_{i \in [m]}\#\{\mathbf{y} \in \{0,1\}^v:
        z_i(\mathbf{y})=t(\mathbf{x})\}$.
    \end{frame}

    \begin{frame}{Running Sum-Check. Or Not?}
        \textbf{Idea \#1:} Run Sum-Check on the sum:
        \begin{equation*}
            \zeta(\mathbf{x}) = \sum_{i \in [m]} \frac{1}{\gamma + z_i(\mathbf{x})} - \frac{\mu(\mathbf{x})}{\gamma+t(\mathbf{x})}
        \end{equation*}

        \vspace{-15px}

        \pause\textcolor{oc-red-8}{\textbf{Problem.}} $\zeta(\mathbf{x})$ is
        a fraction, so we can't quite run the Sum-Check yet. \pause \hspace{0.5px} Idea of logup is to
        split the sum into $\ell$ terms:
        \begin{align*}
            \zeta(\mathbf{x}) &= \textcolor{oc-indigo-7}{\underbrace{\frac{\mu(\mathbf{x})}{\gamma+t(\mathbf{x})} - \frac{1}{\gamma+z_0(\mathbf{x})} - \dots - \frac{1}{\gamma+z_{\ell-2}(\mathbf{x})}}_{\zeta_0(\mathbf{x}), \ell \; \text{terms}}}\\
            & \textcolor{oc-green-8}{\underbrace{- \frac{1}{\gamma+z_{\ell-1}(\mathbf{x})} - \dots - \frac{1}{\gamma + z_{2\ell-2}(\mathbf{x})}}_{\zeta_1(\mathbf{x}), \ell \; \text{terms}}} - \dots 
        \end{align*}

        \vspace{-15px}

        \pause We form $k \approx m/\ell$ \textit{helper columns}
        $\{h_i(\mathbf{x})\}_{i \in [k]}$ that satisfy: (a) $h_i(\mathbf{x})$
        agrees with $\zeta_i(\mathbf{x})$ over $\{0,1\}^v$, (b)
        $\textcolor{oc-indigo-7}{\sum_{\mathbf{x} \in \{0,1\}^v}\sum_{i \in [k]} h_i(\mathbf{x}) = 0}$.
    \end{frame}

    \begin{frame}{Enforcing correct helper columns}
        \textbf{Idea \#2:} Combine $k$ zero-checks using random scalars
        $\{\lambda_i\}_{i \in [k]} \leftarrowS \mathbb{F}$ and merge into a
        single Sum-Check protocol.\pause

        For simplicity, assume each $\zeta_i(\mathbf{x}) = \sum_{j \in
        \mathcal{I}_i} \frac{\mu_j(\mathbf{x})}{\phi_j(\mathbf{x})}$. Note that:
        \begin{equation*}
            \zeta_i(\mathbf{x}) = \sum_{j \in
        \mathcal{I}_i} \frac{\mu_j(\mathbf{x})}{\phi_j(\mathbf{x})} \pause = 
        \frac{\sum_{j \in \mathcal{I}_i}\mu_j(\mathbf{x})\prod_{k \in \mathcal{I}_i \setminus \{j\}}\phi_k(\mathbf{x})}{\prod_{j \in \mathcal{I}_i}\phi_j(\mathbf{x})} \pause = h_i(\mathbf{x})
        \end{equation*}
        
        \vspace{-30px}

        \pause\begin{equation*}
            \implies \textcolor{oc-green-8}{h_i(\mathbf{x})\prod_{j \in \mathcal{I}_i}\phi_j(\mathbf{x}) = \sum_{j \in \mathcal{I}_i}\mu_j(\mathbf{x})\prod_{k \in \mathcal{I}_i \setminus \{j\}}\phi_k(\mathbf{x})}
        \end{equation*}

        \pause\begin{example}
            For $\ell=2$, $h_0(\mathbf{x}) = \frac{\mu(\mathbf{x})}{\gamma+t(\mathbf{x})} - \frac{1}{\gamma+z_0(\mathbf{x})} = \frac{\mu(\mathbf{x})(\gamma+z_0(\mathbf{x})) - \gamma - t(\mathbf{x})}{(\gamma+t(\mathbf{x}))(\gamma+z_0(\mathbf{x}))}$. Thus, we enforce
            the following equality:\pause
            \begin{equation*}
                h_0(\mathbf{x})(\gamma+t(\mathbf{x}))(\gamma+z_0(\mathbf{x})) = \mu(\mathbf{x})(\gamma+z_0(\mathbf{x})) - \gamma - t(\mathbf{x})
            \end{equation*}
        \end{example}
    \end{frame}

    \begin{frame}{Final Sum-Check}
        This way, one runs the Sum-Check on the following function for
        randomly sampled $\boldsymbol{\alpha} \leftarrowS \mathbb{F}^v$:
        \begin{equation*}
            \boxed{\sum_{r \in [k]} \textcolor{oc-indigo-7}{h_r(\mathbf{x})} + \textcolor{oc-green-8}{\mathsf{eq}(\mathbf{x};\boldsymbol{\alpha})\lambda_r\left(h_r(\mathbf{x})\prod_{i \in \mathcal{I}_r}\phi_i(\mathbf{x}) - \sum_{i \in \mathcal{I}_r}\mu_i(\mathbf{x})\prod_{j \in \mathcal{I}_r \setminus \{i\}} \phi_j(\mathbf{x})\right)} = 0}
        \end{equation*}

        \pause\begin{itemize}
            \item The \textcolor{oc-indigo-7}{first term} enforces that
            $\sum_{\mathbf{x} \in \{0,1\}^v}\sum_{r \in [k]}h_r(\mathbf{x}) =
            0$.\pause
            \item The \textcolor{oc-green-8}{second term} checks whether the
            helper columns are consistent with each $\zeta_i(\mathbf{x})$ (see
            prev. slide). For this, we check the equality
            $h_i(\mathbf{x})\prod_{j \in \mathcal{I}_i}\phi_j(\mathbf{x}) = \sum_{j
            \in \mathcal{I}_i}\mu_j(\mathbf{x})\prod_{k \in \mathcal{I}_i \setminus
            \{j\}}\phi_k(\mathbf{x})$ at a random point $\boldsymbol{\alpha}$ by
            multiplying by $\mathsf{eq}(\mathbf{x};\boldsymbol{\alpha})$ and taking a
            linear combination of these equations.\pause
        \end{itemize}

        \begin{remark}
            The larger $\ell$ is (size of each chunk), the more complex
            computations are involved but smaller commitment sizes are required. 
        \end{remark}
    \end{frame}

    \makeclosingframe
\end{document}
